\documentclass[11pt, oneside, article]{memoir}

\linespread{1.15}\selectfont % set leading before the layout, so it is computed on the real \baselineskip
\settrims{0pt}{0pt} % page and stock same size
\settypeblocksize{*}{35.4pc}{*} % {height}{width}{ratio}
\setlrmargins{*}{*}{1} % {spine}{edge}{ratio}
\setulmarginsandblock{0.85in}{0.85in}{*} % top fixed; [lines] grows the bottom margin
\setheadfoot{\onelineskip}{2.66\onelineskip} % {headheight}{footskip}
\setheaderspaces{*}{1.5\onelineskip}{*} % {headdrop}{headsep}{ratio}
\checkandfixthelayout[lines] % round textheight to a whole number of lines

% Running headers: chapter on left, section on right
\makepagestyle{dap}
\makeevenhead{dap}{\small\slshape\leftmark}{}{\small\slshape\rightmark}
\makeoddhead{dap}{\small\slshape\leftmark}{}{\small\slshape\rightmark}
\makeevenfoot{dap}{}{\thepage}{}
\makeoddfoot{dap}{}{\thepage}{}
\makeheadrule{dap}{\textwidth}{\normalrulethickness}
\makepsmarks{dap}{%
  \nouppercaseheads
  \createmark{chapter}{left}{shownumber}{}{.\enspace}%
  \createmark{section}{right}{shownumber}{}{.\enspace}%
}
\pagestyle{dap}

\usepackage{amsthm}
\usepackage{mathtools}

\usepackage[inline]{enumitem}
\usepackage{ifthen}
\usepackage[utf8]{inputenc} %allows non-ascii in bib file
\usepackage{xcolor}

\usepackage[backend=biber, backref=true, maxbibnames = 10, style = alphabetic]{biblatex}
\usepackage[bookmarks=true, colorlinks=true, linkcolor=blue!50!black,
citecolor=orange!50!black, urlcolor=orange!50!black, pdfencoding=unicode]{hyperref}
\usepackage{xurl} % break URLs at any character; avoids biblatex's stretched-out URL spacing
\usepackage[capitalize]{cleveref}

\usepackage{tikz}
\usepackage{pdflscape}
\usepackage{afterpage}

\usepackage{amssymb}
\usepackage{newpxtext}
\usepackage[varg,bigdelims]{newpxmath}
\usepackage{mathrsfs}
\usepackage{dutchcal}
\usepackage[scaled=0.95]{helvet}
\DeclareMathAlphabet{\mathsf}{T1}{phv}{m}{n}
\usepackage{fontawesome}
\usepackage{ebproof}
\usepackage{stmaryrd}

% cleveref %
  \newcommand{\creflastconjunction}{, and\nobreakspace} % serial comma
  \crefformat{enumi}{\card#2#1#3}
  \crefalias{chapter}{section}


% biblatex %
  \addbibresource{Library20260611.bib}

% hyperref %
  \hypersetup{final}

% enumitem %
  \setlist{nosep}
  \setlistdepth{6}


% tikz %

  \usetikzlibrary{
  	cd,
  	math,
  	decorations.markings,
		decorations.pathreplacing,
		decorations.pathmorphing,
  	positioning,
  	arrows.meta,
  	shapes,
		shadows,
		shadings,
  	calc,
  	fit,
  	quotes,
  	intersections,
  	circuits.logic.US,
  }

  \tikzset{
biml/.tip={Glyph[glyph math command=triangleleft, glyph length=.95ex]},
bimr/.tip={Glyph[glyph math command=triangleright, glyph length=.95ex]},
}

\tikzset{
	tick/.style={postaction={
  	decorate,
    decoration={markings, mark=at position 0.5 with
    	{\draw[-] (0,.4ex) -- (0,-.4ex);}}}
  }
}
\newcommand{\tickar}{\begin{tikzcd}[baseline=-0.5ex,cramped,sep=small,ampersand replacement=\&]{}\ar[r,tick]\&{}\end{tikzcd}}
\NewDocumentCommand{\xtickar}{o m}{%
  \begin{tikzcd}[baseline=-0.5ex,cramped,sep=small,ampersand replacement=\&]
    {}\IfNoValueTF{#1}{\ar[r,tick,"{#2}"{above}]}{\ar[r,tick,"{#2}"{above},"{#1}"{below}]}\&{}
  \end{tikzcd}%
}
\tikzset{
	slash/.style={postaction={
  	decorate,
    decoration={markings, mark=at position 0.5 with
    	{\draw[-] (.3ex,.3ex) -- (-.3ex,-.3ex);}}}
  }
}

% oriented WD style (Schultz--Spivak wiring diagrams)
\tikzset{
	oriented WD/.style={
		every to/.style={out=0,in=180,draw},
    label/.style={
    	font=\everymath\expandafter{\the\everymath\scriptstyle},
      inner sep=0pt,
      node distance=2pt and -2pt},
    semithick,
    node distance=1 and 1,
    decoration={markings, mark=at position \stringdecpos with \stringdec},
    ar/.style={postaction={decorate}},
    execute at begin picture={\tikzset{
    	x=\bbx, y=\bby,
      every fit/.style={inner xsep=\bbx, inner ysep=\bby}}}
    },
    string decoration/.store in=\stringdec,
    string decoration={\arrow{stealth}},
    string decoration pos/.store in=\stringdecpos,
    string decoration pos=.7,
    bbx/.store in=\bbx,
    bbx = 1.5cm,
    bby/.store in=\bby,
    bby = 1.5ex,
    bb port sep/.store in=\bbportsep,
    bb port sep=1.5,
    bb port length/.store in=\bbportlen,
    bb port length=4pt,
    bb penetrate/.store in=\bbpenetrate,
    bb penetrate=0,
    bb min width/.store in=\bbminwidth,
    bb min width=1cm,
    bb rounded corners/.store in=\bbcorners,
    bb rounded corners=2pt,
    bb small/.style={
    	bb port sep=1, bb port length=2.5pt, bbx=.4cm, bb min width=.4cm, bby=.7ex},
    bb/.code 2 args={
    	\pgfmathsetlengthmacro{\bbheight}{\bbportsep * (max(#1,#2)+1) * \bby}
      \pgfkeysalso{draw,minimum height=\bbheight,minimum
       width=\bbminwidth,outer sep=0pt,
         rounded corners=\bbcorners,thick,
         prefix after command={\pgfextra{\let\fixname\tikzlastnode}},
         append after command={\pgfextra{\draw
            \ifnum #1=0{} \else foreach \i in {1,...,#1} {
            	($(\fixname.north west)!{\i/(#1+1)}!(\fixname.south west)$) +(-\bbportlen,0) coordinate (\fixname_in\i) -- +(\bbpenetrate,0) coordinate (\fixname_in\i')}\fi
            \ifnum #2=0{} \else foreach \i in {1,...,#2} {
            	($(\fixname.north east)!{\i/(#2+1)}!(\fixname.south east)$) +(-\bbpenetrate,0) coordinate (\fixname_out\i') -- +(\bbportlen,0) coordinate (\fixname_out\i)}\fi;
           }}}
		},
		bb name/.style={
     	append after command={
				\pgfextra{\node[anchor=north] at (\fixname.north) {#1};}
			}
		}
}

% a single {1}{1} box with an input and an output wire;
% #1 = input-wire label, #2 = output-wire label (either may be empty)
\newcommand{\boxpic}[2]{%
  \begin{tikzpicture}[oriented WD, bb small, baseline=(M.south)]
    \node[bb={1}{1}] (M) {};
    \node[left=-1pt] at (M_in1) {\scriptsize #1};
    \node[right=-1pt] at (M_out1) {\scriptsize #2};
  \end{tikzpicture}%
}

% a single {0}{1} box with only an output wire; #1 = output-wire label
\newcommand{\boxpicout}[1]{%
  \begin{tikzpicture}[oriented WD, bb small, baseline=(M.south)]
    \node[bb={0}{1}] (M) {};
    \node[right=-1pt] at (M_out1) {\scriptsize #1};
  \end{tikzpicture}%
}


  % amsthm + aliascnt (so cleveref prints the correct environment name) %
\usepackage{aliascnt}

\theoremstyle{definition}
\newtheorem{definitionx}{Definition}[section]
\crefname{definitionx}{Definition}{Definitions}

\newaliascnt{examplex}{definitionx}
\newtheorem{examplex}[examplex]{Example}
\aliascntresetthe{examplex}
\crefname{examplex}{Example}{Examples}

\newaliascnt{warning}{definitionx}
\newtheorem{warning}[warning]{Warning}
\aliascntresetthe{warning}
\crefname{warning}{Warning}{Warnings}

\newaliascnt{remarkx}{definitionx}
\newtheorem{remarkx}[remarkx]{Remark}
\aliascntresetthe{remarkx}
\crefname{remarkx}{Remark}{Remarks}

\newaliascnt{notation}{definitionx}
\newtheorem{notation}[notation]{Notation}
\aliascntresetthe{notation}
\crefname{notation}{Notation}{Notations}

\newaliascnt{conventionx}{definitionx}
\newtheorem{conventionx}[conventionx]{Convention}
\aliascntresetthe{conventionx}
\crefname{conventionx}{Convention}{Conventions}

\theoremstyle{plain}

\newaliascnt{theorem}{definitionx}
\newtheorem{theorem}[theorem]{Theorem}
\aliascntresetthe{theorem}
\crefname{theorem}{Theorem}{Theorems}

\newaliascnt{proposition}{definitionx}
\newtheorem{proposition}[proposition]{Proposition}
\aliascntresetthe{proposition}
\crefname{proposition}{Proposition}{Propositions}

\newaliascnt{corollary}{definitionx}
\newtheorem{corollary}[corollary]{Corollary}
\aliascntresetthe{corollary}
\crefname{corollary}{Corollary}{Corollaries}

\newaliascnt{lemma}{definitionx}
\newtheorem{lemma}[lemma]{Lemma}
\aliascntresetthe{lemma}
\crefname{lemma}{Lemma}{Lemmas}

\newaliascnt{fact}{definitionx}
\newtheorem{fact}[fact]{Fact}
\aliascntresetthe{fact}
\crefname{fact}{Fact}{Facts}

\newaliascnt{conjecture}{definitionx}
\newtheorem{conjecture}[conjecture]{Conjecture}
\aliascntresetthe{conjecture}
\crefname{conjecture}{Conjecture}{Conjectures}

\newtheorem*{theorem*}{Theorem}
\newtheorem*{proposition*}{Proposition}
\newtheorem*{corollary*}{Corollary}
\newtheorem*{lemma*}{Lemma}
\newtheorem*{warning*}{Warning}

\newenvironment{example}
  {\pushQED{\qed}\renewcommand{\qedsymbol}{$\lozenge$}\examplex}
  {\popQED\endexamplex}

 \newenvironment{remark}
  {\pushQED{\qed}\renewcommand{\qedsymbol}{$\lozenge$}\remarkx}
  {\popQED\endremarkx}

  \newenvironment{definition}
  {\pushQED{\qed}\renewcommand{\qedsymbol}{$\lozenge$}\definitionx}
  {\popQED\enddefinitionx}

  \newenvironment{convention}
  {\pushQED{\qed}\renewcommand{\qedsymbol}{$\lozenge$}\conventionx}
  {\popQED\endconventionx}


%--------Single symbols--------%

\DeclareSymbolFont{stmry}{U}{stmry}{m}{n}
\DeclareMathSymbol\fatsemi\mathop{stmry}{"23}

\DeclareFontFamily{U}{mathx}{\hyphenchar\font45}
\DeclareFontShape{U}{mathx}{m}{n}{
      <5> <6> <7> <8> <9> <10>
      <10.95> <12> <14.4> <17.28> <20.74> <24.88>
      mathx10
      }{}
\DeclareSymbolFont{mathx}{U}{mathx}{m}{n}
\DeclareFontSubstitution{U}{mathx}{m}{n}
\DeclareMathAccent{\widecheck}{0}{mathx}{"71}


%--------Renewed commands--------%

\renewcommand{\ss}{\subseteq}

%--------Other Macros--------%


\DeclarePairedDelimiter{\present}{\langle}{\rangle}
\DeclarePairedDelimiter{\pairing}{\langle}{\rangle}
\DeclarePairedDelimiter{\copair}{[}{]}
\DeclarePairedDelimiter{\floor}{\lfloor}{\rfloor}
\DeclarePairedDelimiter{\ceil}{\lceil}{\rceil}
\DeclarePairedDelimiter{\corners}{\ulcorner}{\urcorner}
\DeclarePairedDelimiter{\ihom}{[}{]}
\DeclarePairedDelimiter{\absval}{\lvert}{\rvert}

\DeclareMathOperator{\Hom}{Hom}
\DeclareMathOperator{\Aut}{Aut}
\DeclareMathOperator{\End}{End}
\DeclareMathOperator{\Iso}{Iso}
\DeclareMathOperator{\Mor}{Mor}
\DeclareMathOperator{\dom}{dom}
\DeclareMathOperator{\cod}{cod}
\DeclareMathOperator{\idy}{idy}
\DeclareMathOperator*{\colim}{colim}
\DeclareMathOperator{\im}{im}
\DeclareMathOperator{\ob}{Ob}
\DeclareMathOperator{\Tr}{Tr}
\DeclareMathOperator{\el}{El}


\newcommand{\cat}[1]{\mathcal{#1}}%a generic category
\newcommand{\Cat}[1]{\mathbf{#1}}%a named category
\newcommand{\fun}[1]{\mathrm{#1}}%a function
\newcommand{\Fun}[1]{\text{\normalfont\textsf{#1}}}%a named functor
%\arr is registered via \trackTerm below (after the term-tracking machinery is defined)
\providecommand{\Data}{\Cat{Data}}%category of adaptive-arrangement data (with integrator)
\providecommand{\SMC}{\Cat{SMC}}%symmetric monoidal categories (codomain slice)

\newcommand{\notsafe}[1]{{#1}}
\makeatletter
\NewDocumentCommand{\trackTermSymbol}{m m}{%
  \expandafter\newcommand\csname termraw@#1\endcsname{\notsafe{#2}}%
}
\NewDocumentCommand{\termraw}{m}{%
  \@ifundefined{termraw@#1}%
    {\PackageError{termraw}{Term '#1' not registered}{Use \string\trackTermSymbol{#1}{...} first.}}%
    {\@nameuse{termraw@#1}}%
}
\DeclareRobustCommand{\termref}[2]{\notsafe{\hypersetup{linkcolor=black}\hyperlink{term.#1}{#2}}}
\NewDocumentCommand{\trackTerm}{m m m}{%
  \trackTermSymbol{#2}{#3}%
  \newcommand{#1}{\termref{#2}{\termraw{#2}}}%
}
\NewDocumentCommand{\defineTerm}{m}{\notsafe{\hypertarget{term.#1}{\termraw{#1}}}}
\makeatother
% To disable semantic notation links, replace the \termref line above with:
% \DeclareRobustCommand{\termref}[2]{\notsafe{#2}}


\newcommand{\id}{\mathrm{id}}
\newcommand{\then}{\mathbin{\fatsemi}}

\newcommand{\too}{\longrightarrow}
\newcommand{\To}[2][]{\xrightarrow[#1]{#2}}
\renewcommand{\Mapsto}[1]{\xmapsto{#1}}
\newcommand{\from}{\leftarrow}
\newcommand{\From}[1]{\xleftarrow{#1}}
\newcommand{\inj}{\hookrightarrow}
\newcommand{\surj}{\twoheadrightarrow}
\newcommand{\tto}{\rightrightarrows}
\renewcommand{\iff}{\Leftrightarrow}
\newcommand{\card}{\,^{\#}}
\makeatletter
\newcommand{\leftvec}[1]{\mathpalette\leftvec@{#1}}
\newcommand{\leftvec@}[2]{\reflectbox{$\m@th#1\vec{\reflectbox{$\m@th#1#2$}}$}}
\makeatother
\newcommand{\bk}[2]{#1_{\leftvec{#2}}}


\newcommand{\inv}{^{-1}}
\newcommand{\op}{^\tn{op}}

\newcommand{\tn}[1]{\textnormal{#1}}

\newcommand{\nn}{\mathbb{N}}
\newcommand{\rr}{\mathbb{R}}


\trackTerm{\finset}{finset}{\Cat{Fin}}
\trackTerm{\smset}{smset}{\Cat{Set}}
\trackTerm{\smsetiso}{smsetiso}{\Cat{Set}_\cong}
\trackTerm{\smcat}{smcat}{\Cat{Cat}}
\trackTerm{\vect}{vect}{\Cat{Vect}}
\trackTerm{\vectiso}{vectiso}{\Cat{Vect}_\cong}
\trackTerm{\sbf}{sbf}{\mathit{SBF}}

\newcommand{\List}{\Fun{List}}


\trackTerm{\yon}{yon}{\mathcal{y}}
\trackTerm{\poly}{poly}{\Cat{Poly}}

% lenses
\newcommand{\biglens}[2]{
     \begin{bmatrix}{\vphantom{f_f^f}#2} \\ {\vphantom{f_f^f}#1} \end{bmatrix}
}
\newcommand{\littlelens}[2]{
     \begin{bsmallmatrix}{\vphantom{f}#2} \\ {\vphantom{f}#1} \end{bsmallmatrix}
}
\newcommand{\lens}[2]{
  \relax\if@display
     \biglens{#1}{#2}
  \else
     \littlelens{#1}{#2}
  \fi
}


\newcommand{\qand}{\quad\text{and}\quad}
\newcommand{\qqand}{\qquad\text{and}\qquad}
\newcommand{\qqresp}{,\qquad\text{resp.}\quad}

\newcommand{\coalg}{\tn{-}\Cat{Coalg}}
\newcommand{\alg}{\tn{-}\Cat{Alg}}
\newcommand{\Coalg}{\Fun{Coalg}}
\newcommand{\enrcat}[1]{#1\tn{-}\smcat}% category of #1-enriched categories

\trackTerm{\pc}{pc}{\mathbb{P}\Cat{C}}
\trackTerm{\pciso}{pciso}{\mathbb{P}\Cat{C}_{\cong}}
\trackTerm{\arr}{rwd}{\mathbb{A}\Cat{rr}}% abstract adaptive-arrangement operad

%--------Paper-specific macros--------%

\trackTerm{\mfd}{mfd}{\Cat{Mfd}}
\trackTerm{\Store}{Store}{\Fun{Store}}
\newcommand{\inc}{\Fun{inc}}
\trackTermSymbol{cot}{\fun{cot}}
\RenewDocumentCommand{\cot}{}{\termref{cot}{\termraw{cot}}}
\NewDocumentCommand{\cotof}{m}{\termref{cot}{\termraw{cot}}(#1)}
\trackTermSymbol{chiQ}{\chi}
\newcommand{\chiQ}[1]{\termref{chiQ}{\termraw{chiQ}}_{#1}}
\trackTermSymbol{Psisem}{\Psi}
\newcommand{\Psisem}[1]{\termref{Psisem}{\termraw{Psisem}}_{#1}}
\newcommand{\ev}{\fun{ev}}
\newcommand{\coev}{\fun{coev}}
\newcommand{\src}{\mathrm{src}}
\newcommand{\tgt}{\mathrm{tgt}}
\newcommand{\Part}{\bullet}
\newcommand{\boxob}{\mathrm{box}}
\trackTerm{\bang}{bang}{\mathord{!}}
\newcommand{\tpow}[1]{^{\otimes #1}}
\providecommand{\Sm}{\mathrm{Sm}}%the smooth adaptive-arrangement datum (instance of def.rewiring_datum); \sarr=\arr_{\Sm}
\trackTerm{\sarr}{srwd}{\mathbb{A}\Cat{rr}_{\Sm}}
\trackTerm{\Para}{Para}{\mathbb{P}\Cat{ara}}
\newcommand{\para}[3][]{\Para_{#2}^{#1}(#3)}
\trackTerm{\rmfdc}{rmfdc}{\Cat{RMfdC}}
\trackTerm{\rvect}{rvect}{\Cat{RVect}}
\trackTerm{\SLens}{SLens}{\Cat{SLens}}% submersion lenses (def. in proof of prop.euler_submersion_lenses)
\newcommand{\rv}[1]{\mathbf{#1}}% a reactive vector space (object of \rvect): \rv Q=(Q,\sharpR_Q), carrier Q
\trackTermSymbol{Lens}{\Cat{Lens}}
\newcommand{\Lens}[1]{\termref{Lens}{\termraw{Lens}}_{#1}}
\newcommand{\lmfd}{\Lens{\mfd}}
\trackTermSymbol{LensFun}{\Cat{Lens}}% the backward comonad on \Lens{\cat C} induced by a monad (def. at prop.backward_comonad)
\newcommand{\LensFun}[1]{\termref{LensFun}{\termraw{LensFun}}^{#1}}
\trackTermSymbol{LensMor}{\Cat{Lens}}% the induced coKleisli functor from a pair (F,\theta) (def. at lem.lens_T_functoriality)
\newcommand{\LensMor}[2]{\termref{LensMor}{\termraw{LensMor}}_{#1}^{#2}}
\trackTermSymbol{Lcokl}{\Cat{Lens}_{\cat C}^{\Fun T}}
\newcommand{\Lcokl}[2]{\termref{Lcokl}{\Cat{Lens}_{#1}^{#2}}}
\newcommand{\plmfd}{\Lcokl{\mfd}{\rr}}
\trackTermSymbol{prokl}{\cat{C}_{T}}% Kleisli category of a promonad T on \cat{C} (def. at prop.promonad_io_functor)
\newcommand{\prokl}[2]{\termref{prokl}{{#1}_{#2}}}% Kleisli category of the promonad #2 on the category #1
\newcommand{\CT}{\prokl{\cat{C}}{T}}% Kleisli category of a generic promonad T on \cat{C}
\newcommand{\ltens}{\boxtimes}% monoidal product on a lens category (dedicated, base-independent)
\newcommand{\bigltens}{\mathop{\vcenter{\hbox{\begin{tikzpicture}[line width=0.6pt,line join=miter]\def\s{5.5pt}\draw(-\s,-\s)rectangle(\s,\s);\draw(-\s,-\s)--(\s,\s);\draw(-\s,\s)--(\s,-\s);\end{tikzpicture}}}}\displaylimits}% indexed lens tensor (big \ltens), drawn so line weight is independent of size
\newcommand{\ltpow}[1]{^{\ltens #1}}% lens-tensor power (cf. \tpow for base/poly tensor power)
\newcommand{\inpt}[1]{#1^{-}}
\newcommand{\outp}[1]{#1^{+}}
\newcommand{\lensob}[1]{\binom{\inpt{#1}}{\outp{#1}}}% vertical lens-object notation; to revert, change to (\inpt{#1},\outp{#1})
\newcommand{\lensobs}[2]{\binom{\inpt{#1}\otimes\inpt{#2}}{\outp{#1}\otimes\outp{#2}}}% monoidal product of two lens objects
\trackTerm{\potd}{potd}{\fun{d}}
\trackTerm{\potalg}{potalg}{\mathbf z}% T-monoid / potential algebra: \potalg=(z,e_z,m_z,\alpha), carrier z (defined at def.T_monoid)
\newcommand{\md}{\mathsf{Md}}% finite set of modes; element i\in\md (prop.moding, prop.moded_algebra)
\newcommand{\conf}{\mathrm{conf}}
\newcommand{\phase}{\mathrm{phase}}
\newcommand{\lrn}{\mathrm{lrn}}
\newcommand{\euc}{\mathrm{Euc}}

\newcommand{\Phip}[1]{\Phi'_{#1}}% polynomial interpretation; subscript is always the interpretation symbol (\interp=\mathfrak p), abbreviating the full 4-tuple datum. Never a partial unfolding like \Phi'_{\cot,(\yon,\potd)}.
\trackTermSymbol{Phiphase}{\Phi_{\phase}}
\newcommand{\Phiphase}{\termref{Phiphase}{\termraw{Phiphase}}}
\trackTermSymbol{Phiconf}{\Phi_{\conf}}
\newcommand{\Phiconf}{\termref{Phiconf}{\termraw{Phiconf}}}
\trackTermSymbol{sharpEuc}{\sharp_{\euc}}
\newcommand{\sharpEuc}{\termref{sharpEuc}{\termraw{sharpEuc}}}
\trackTermSymbol{dyn}{\mathrm{dyn}}
\newcommand{\dyn}{\termref{dyn}{\termraw{dyn}}}
\trackTermSymbol{dynphase}{\mathrm{dyn}_{\mathrm{phase}}}% phase-integrator dynamics (def. at eqn.phase_integrator)
\newcommand{\dynphase}{\termref{dynphase}{\termraw{dynphase}}}
\trackTerm{\intg}{intg}{\mathfrak{i}}
\trackTerm{\interp}{interp}{\mathfrak{p}}
\trackTerm{\interpsm}{interpsm}{\mathrm{sm}}% smooth (cotangent) polynomial interpretation datum; concrete named choice (roman, like \conf/\phase), not the generic \interp=\mathfrak p of ch.framework
\trackTermSymbol{sharpR}{\sharp}
\newcommand{\sharpR}{\termref{sharp_map}{\termraw{sharpR}}}
\trackTermSymbol{sharp_symp}{\sharp}
\trackTermSymbol{sharpS}{\sharp}
\newcommand{\sharpS}{\termref{sharp_symp}{\termraw{sharpS}}}

\font\MnSyCfont=MnSymbolC10
\font\MnSyCfonts=MnSymbolC7
\font\MnSyCfontss=MnSymbolC5
\newcommand{\smallsquare}{\mathchoice
  {\mbox{\MnSyCfont\symbol{"68}}}%
  {\mbox{\MnSyCfont\symbol{"68}}}%
  {\mbox{\MnSyCfonts\symbol{"68}}}%
  {\mbox{\MnSyCfontss\symbol{"68}}}}
\newcommand{\blank}{\notsafe{\mathord{\mkern1mu\smallsquare\mkern2mu}}}

\newcommand{\thanksAFOSR}[1]{This material is based upon work supported by the Air Force Office of Scientific Research under award number #1}

\newcommand{\dnote}[1]{\quad{\color{red!80!black}$\lozenge$David says:\quad#1\quad$\lozenge$}}
\newcommand{\mnote}[1]{\quad{\color{blue!80!black}$\lozenge$Maxi says:\quad#1\quad$\lozenge$}}


%----Changeable document parameters----%

% \linespread is set at the top, before \checkandfixthelayout, so it also drives the page layout
%\allowdisplaybreaks
\setsecnumdepth{subsection}
\settocdepth{section}
\setlength{\parindent}{15pt}
\setcounter{tocdepth}{1}


%---------------Document---------------%
\begin{document}

\title{Compositional Linearization for Superposition and Hamiltonianization}

\author{
Maximilien P\'eroux
\and
David I. Spivak}

\date{\vspace{-.2in}}

\maketitle

\begin{abstract}
test
\end{abstract}

\begin{itemize}
  \item Def of monoidal promonad. Main example is quadratic form.
  \item Any generaly theory of quadratic forms can be translated?
  \item Algebra for monoidal promonads.
  \item Moore internalization for promondic lenses. Check if true.
  \item Lens for monoidal promonads. 
  \item Quick reminder about Para.
  \item Monoidal product on Lin. 
  \item Check moidal retraction/section between Lin and Crit.
  \item Lift to Coalg over Poly.
  \item Superposition and physics. Lots of examples.
  \item Some relation to Poincaré point of view. 
  \item Any relation to linearization of system of differential equations? PDEs?
\end{itemize}

%\tableofcontents*

%====================================================================
\chapter{Introduction}\label{ch.intro}
%====================================================================


%====================================================================
\chapter{Promonads}\label{ch.promonas}
%====================================================================

Say that a functor is \emph{io} if it is identity on objects. 

\begin{proposition}\label{prop.promonad_io_functor}
Let $\cat{C}$ be a category. The following sets are in bijection:
\begin{enumerate}
	\item a monad $(T,\eta,\mu)$ in the category of profunctors $\cat{C}\tickar\cat{C}$;
	\item a category $\defineTerm{prokl}$ equipped with an io functor $\cat{C}\To{\eta}\CT$.
\end{enumerate}
We call these \emph{promonads} on $\cat{C}$, and we call $\CT$ the \emph{Kleisli category} of the promonad $T$; the underlying profunctor of $T$ is recovered as $T(x,y)=\CT(x,y)$. Let $T,T'$ be promonads on $\cat{C}$. The following sets are in bijection
\begin{enumerate}
	\item monad morphisms $T\to T'$;
	\item functors $\CT\to\prokl{\cat{C}}{T'}$ commuting with $\eta$ and $\eta'$.
\end{enumerate}
\end{proposition}
\begin{proof}
**
\end{proof}

\begin{proposition}\label{prop.kleisli_cokleisli_promonad}
Let $\cat{C}$ be a category and $T$ a monad (resp.\ comonad) on it. Then the io functor $\cat{C}\to\CT$ into the Kleisli (resp.\ coKleisli) category of $T$ is a promonad on $\cat{C}$, which we again denote $T$; its underlying profunctor is $T(x,y)=\cat{C}(x,Ty)$ (resp.\ $T(x,y)=\cat{C}(Tx,y)$). In particular, the notation $\CT$ from \cref{prop.promonad_io_functor} is unambiguous.
\end{proposition}
\begin{proof}
**
\end{proof}

\begin{proposition}\label{prop.grothendieck_monoid_promonad}
Let $\cat{C}$ be a category and $M\colon\cat{C}\op\to\Cat{Mon}$ a functor to the category of set-theoretic monoids, i.e.\ categories with one object. Its Grothendieck construction $\int_\cat{C}M$ is the Kleisli category of a promonad on $\cat{C}$.
\end{proposition}

\begin{example}\label{ex.symmetric_bilinear_forms}
Given a vector space $V$, there is a vector space of symmetric bilinear forms $B\colon V\otimes V\to\rr$, and these pull back along linear maps. Taking the underlying additive monoid, we obtain a functor $\sbf\colon\vect\op\to\Cat{Mon}$. Its Grothendieck construction has vector spaces as objects $\ob\int_\vect\sbf=\ob\vect$ and a map $V\to V'$ consists of a pair $(f,B)$, where $f\colon V\to V'$ is a linear map and $B$ is a symmetric bilinear form on $V$. The identity is $(\id_V,0)$ and the composite $(f,B)\then(f',B')$ is $(f\then f', B+((f\otimes f)\then B'))$. By \cref{prop.grothendieck_monoid_promonad} this is the Kleisli category of a promonad on $\vect$.
\end{example}

\begin{definition}\label{def.promonad_algebra_coalgebra}
An \emph{algebra} (resp.\ coalgebra) for a promonad $T$ on $\cat{C}$ consists of a pair $(C,\alpha)$ where $C\in\cat{C}$ is an object and
\[
\alpha\colon T(\blank,C)\to\cat{C}(\blank,C)
\qqresp
\alpha\colon T(C,\blank)\to\cat{C}(C,\blank)
\]
is a natural transformation making the following diagrams commute\dnote{Make sense of lefthand map in righthand diagram}
\[
\begin{tikzcd}
  \cat{C}(\blank, C)\ar[r, "\eta"]\ar[d,equal]&
  T(\blank, C)\ar[d, "\alpha"]\\
  \cat{C}(\blank, C)\ar[r, equal]&
  \cat{C}(\blank, C)
\end{tikzcd}
\hspace{.5in}
\begin{tikzcd}
	(T\circ T)(\blank, C)\ar[r, "\mu"]\ar[d, "T\circ\alpha"']&
	T(\blank, C)\ar[d, "\alpha"]\\
	T(\blank, C)\ar[r, "\alpha"']&
	\cat{C}(\blank, C)
\end{tikzcd}
\]
respectively (for coalgebras)
\[
\begin{tikzcd}
  \cat{C}(C, \blank)\ar[r, "\eta"]\ar[d,equal]&
  T(C, \blank)\ar[d, "\alpha"]\\
  \cat{C}(C, \blank)\ar[r, equal]&
  \cat{C}(C, \blank)
\end{tikzcd}
\hspace{.5in}
\begin{tikzcd}
	(T\circ T)(C, \blank)\ar[r, "\mu"]\ar[d, "\alpha\circ T"']&
	T(C, \blank)\ar[d, "\alpha"]\\
	T(C, \blank)\ar[r, "\alpha"']&
	\cat{C}(C, \blank)
\end{tikzcd}
\]
\end{definition}

\begin{proposition}\label{prop.kleisli_promonad_algebras}
Let $T$ be a monad (resp.\ comonad) on $\cat{C}$, considered as a promonad as in \cref{prop.kleisli_cokleisli_promonad}. Then there is an isomorphism between the category of monad-algebras for $T$ and the category of promonad-algebras for $T$ (resp.\ between the category of comonad-coalgebras and promonad-coalgebras).
\end{proposition}


\begin{proposition}\label{prop.monoidal_promonad_io_functor}
Let $(\cat{C},I,\otimes)$ be a monoidal category. The following are equivalent:
\begin{enumerate}
	\item promonads $(T,\eta,\mu)$ for which $T\colon\cat{C}\op\times\cat{C}\to\smset$ is lax monoidal and $\eta$ and $\mu$ are monoidal natural transformations;
	\item monoidal structures $(J,\odot)$ on the Kleisli category $\CT$ for which the io functor $\cat{C}\To{\eta}\CT$ is strong monoidal.
\end{enumerate}
We call these \emph{monoidal promonads} on $\cat{C}$.
\end{proposition}
\begin{proof}
**
\end{proof}

%--------------------------------------------------------------------
\section{Promonads and lenses}\label{sec.promonads_lenses}
%--------------------------------------------------------------------

\subsection{The category of lenses}\label{sec.category_of_lenses}
We first recall the category $\defineTerm{Lens}_{\cat{C}}$ of lenses in a symmetric monoidal category $(\cat{C},I,\otimes)$; see \cite{hedges2017coherence} and \cite[Def.~2.9]{spivak2019generalized}. Let $(\cat{C}, I,\otimes)$ be a symmetric monoidal category.

An \emph{interface} in $\cat C$ is a pair $\lensob c$ whose output part $\outp c$ carries a comonoid structure $(\outp c,\varepsilon,\delta)$.%
\footnote{
When $(\cat{C},1,\times)$ is cartesian monoidal, every object has a unique comonoid structure and every morphism is automatically a comonoid homomorphism~\cite{fox1976coalgebras}.
}
A \emph{lens} $f\colon\lensob c\to\lensob d$ consists of a forward map $\outp f\colon\outp c\to\outp d$ (a comonoid homomorphism) and a backward map $\inpt f\colon\outp c\otimes\inpt d\to\inpt c$ (any map in $\cat{C}$). Such a morphism can be depicted as follows:
\begin{equation}\label{eqn.lens_map}
\begin{tikzpicture}[oriented WD, bb port length=5pt, bb port sep=1, bb min width=.4cm, bbx=.5cm, bby=.7ex, baseline=(fo)]
  \node[bb={1}{1}] (fo) {$\outp f$};
  \node[bb={1}{2}, above left=6 and 1.5 of fo] (fi) {$\inpt f$};
  \node[circle, draw, fill=black, inner sep=0, minimum size=3pt, left=1 of fo] (dot) {};
  \coordinate (co) at ($(dot)+(-3,0)$);
  \coordinate (do) at ($(fo.east)+(1,0)$);
  \coordinate (ci) at (co |- fi_in1);
  \coordinate (di) at (do |- fi_out2);
  \node[left=0 of co] {$\outp c$};
  \node[right=0 of do] {$\outp d$};
  \node[left=0 of ci] {$\inpt c$};
  \node[right=0 of di] {$\inpt d$};
  \draw[ar] (co) to (dot);
  \draw[ar] (dot) to (fo_in1);
  \draw[ar] (fo_out1) to (do);
  \draw[ar] (fi_in1) to (ci);
  \draw[ar] (di) to (fi_out2);
  \draw[ar] (dot) to[out=30, in=0] (fi_out1);
\end{tikzpicture}
\end{equation}
The identity lens has $\outp{\id_{\lensob{c}}}\coloneqq\id_{\outp{c}}$ and $\inpt{\id_{\lensob{c}}}\coloneqq\varepsilon\otimes\id_{\inpt{c}}$, i.e.\ use the counit to discard the $\outp c$ part.

Lenses compose as follows: given also $g\colon \lensob{d}\to\lensob{e}$, the composite $f\then g$ is given by
\[
\outp{(f\then g)}=\outp f\then\outp g
\qqand
\inpt{(f\then g)}=
	(\delta_{\outp{c}}\otimes\inpt{e})\then
	(\outp{c}\otimes\outp{f}\otimes\inpt{e})\then
	(\outp{c}\otimes\inpt{g})\then
	\inpt{f}.
\]
This dense notation can be depicted in a string-diagram as follows:
\begin{equation*}
\begin{tikzpicture}[oriented WD, bb port length=5pt, bb port sep=1, bb min width=.4cm, bbx=.5cm, bby=.7ex, baseline=(fo)]
  \node[bb={1}{1}] (fo) {$\outp f$};
  \node[bb={1}{1}, right=5 of fo] (go) {$\outp g$\vphantom{$\inpt f$}};
  \node[bb={1}{2}, above left=6 and 1.5 of fo] (fi) {$\inpt f$};
  \node[bb={1}{2}, above left=6 and 1.8 of go] (gi) {$\inpt g$\vphantom{$\inpt f$}};
  \node[circle, draw, fill=black, inner sep=0, minimum size=3pt, left=1 of fo] (dot1) {};
  \node[circle, draw, fill=black, inner sep=0, minimum size=3pt] at ($(fo.east)!.75!(go.west)$) (dot2) {};
  \coordinate (co) at ($(dot1)+(-3,0)$);
  \coordinate (eo) at ($(go.east)+(1,0)$);
  \coordinate (ci) at (co |- fi_in1);
  \coordinate (ei) at (eo |- gi_out2);
  \node[left=0 of co] {$\outp c$};
  \node[right=0 of eo] {$\outp e$};
  \node[left=0 of ci] {$\inpt c$};
  \node[right=0 of ei] {$\inpt e$};
  \draw[ar] (co) to (dot1);
  \draw[ar] (dot1) to (fo_in1);
  \draw[ar] (fo_out1) to node[above, font=\scriptsize] {$\outp d$} (dot2);
  \draw[ar] (dot2) to (go_in1);
  \draw[ar] (go_out1) to (eo);
  \draw[ar] (fi_in1) to (ci);
  \draw[ar] (ei) to (gi_out2);
  \draw[ar] (gi_in1) to[out=180, in=0] node[above, font=\scriptsize] {$\inpt d$} (fi_out2);
  \draw[ar] (dot1) to[out=30, in=0] (fi_out1);
  \draw[ar] (dot2) to[out=30, in=0] (gi_out1);
\end{tikzpicture}
\end{equation*}

\begin{definition}\label{def.lens}
For a symmetric monoidal category $(\cat C, I,\otimes)$, the category $\Lens{\cat C}$ of \emph{interfaces} and \emph{lenses} has as objects the pairs $\lensob c$ with $\outp c$ a comonoid $(\varepsilon_{\outp c},\delta_{\outp c})$, and as morphisms $\lensob c\to\lensob d$ the pairs $\binom{\inpt f}{\outp f}$ with $\outp f\colon\outp c\to\outp d$ a comonoid homomorphism and $\inpt f\colon\outp c\otimes\inpt d\to\inpt c$ any map in $\cat{C}$, with identity and composition as above.

It is symmetric monoidal with unit $\binom{I}{I}$ and
\[
\lensob{c}\ltens\lensob{d}\coloneqq\lensobs{c}{d}.
\qedhere
\]
\end{definition}

\begin{proposition}\label{prop.monoidal_promonad_lenses}
Let $\cat{C}$ be a symmetric monoidal category and $T$ a monoidal promonad on it. Then there is a promonad on $\Lens{\cat{C}}$ whose Kleisli category, denoted $\Lcokl{\cat{C}}{T}$, has the same objects as $\Lens{\cat{C}}$, and in which a map $\lensob{c}\to\lensob{d}$ consists of a comonoid homomorphism $\outp c\to\outp d$ in $\cat{C}$ and a map $\outp c\otimes\inpt d\to\inpt c$ in $\CT$.
\end{proposition}



\printbibliography

\end{document}
